\heading{热学-免修考-22秋}
\noteinfo{题目来源于上传的 OneNote 扫描件。扫描件只保留了部分题目图片和手写提纲；以下按可辨认内容整理，并在答案中注明无法唯一确定的几何量。}

\textbf{一、简答题（扫描件可辨认部分）}

\begin{question}
在热物理学中，根据组成系统的粒子的内禀属性和遵守的统计规律，微观粒子可以分为哪几类？写出这些系统中的微观粒子按能量的分布规律的表达式，并说明 Maxwell 速率分布为什么可看作平衡态经典理想气体的实际分布律。
\begin{solution}
可分为经典粒子、玻色子和费米子。能级 $\varepsilon_i$、简并度 $g_i$ 的平均占有数为
\[
\bar n_i^{\rm MB}=g_i e^{(\mu-\varepsilon_i)/(k_BT)},\qquad
\bar n_i^{\rm BE}=\frac{g_i}{e^{(\varepsilon_i-\mu)/(k_BT)}-1},\qquad
\bar n_i^{\rm FD}=\frac{g_i}{e^{(\varepsilon_i-\mu)/(k_BT)}+1}.
\]
经典稀薄极限下，速度分布为
\[
f(\mathbf v)=n\left(\frac{m}{2\pi k_BT}\right)^{3/2}e^{-mv^2/(2k_BT)},
\]
速率分布为
\[
f(v)=4\pi n\left(\frac{m}{2\pi k_BT}\right)^{3/2}v^2e^{-mv^2/(2k_BT)}.
\]
由于宏观系统 $N$ 极大，最概然分布附近相对涨落为 $O(N^{-1/2})$，故实际观测分布与最概然分布一致。
\end{solution}
\end{question}

\begin{question}
一定量的气体先经过等体过程使温度升高一倍，再经等温过程使体积膨胀到原来的两倍。求末态的黏度 $\eta$、热导率 $\kappa$、扩散系数 $D$ 和平均自由程 $\bar\lambda$ 分别为初态的多少倍。
\begin{solution}
末态 $T=2T_0$、$n=n_0/2$。稀薄气体中
\[
\eta\propto\sqrt T,
\qquad \kappa\propto\sqrt T,
\qquad \bar\lambda\propto n^{-1},
\qquad D\propto \bar\lambda\sqrt T.
\]
所以
\[
\frac{\eta}{\eta_0}=\sqrt2,
\qquad
\frac{\kappa}{\kappa_0}=\sqrt2,
\qquad
\frac{D}{D_0}=2\sqrt2,
\qquad
\frac{\bar\lambda}{\bar\lambda_0}=2.
\]
\end{solution}
\end{question}

\begin{question}
实验上测得某种气体在压强为 $p$、体积为 $V$、温度为 $T$ 的状态附近，体膨胀系数和等温压缩系数满足
\[
\alpha=\frac{n\nu a}{T^{n+1}V}+\frac{\nu R}{pV},
\qquad
\kappa_T=\frac{\nu RT}{p^2V},
\]
其中 $n,\nu,a,R$ 为正的常量。试确定该气体系统的状态方程。
\begin{solution}
由定义
\[
\alpha=\frac1V\left(\frac{\partial V}{\partial T}\right)_p,
\qquad
\kappa_T=-\frac1V\left(\frac{\partial V}{\partial p}\right)_T,
\]
得
\[
\left(\frac{\partial V}{\partial T}\right)_p=\frac{n\nu a}{T^{n+1}}+\frac{\nu R}{p},
\qquad
\left(\frac{\partial V}{\partial p}\right)_T=-\frac{\nu RT}{p^2}.
\]
对 $T$ 积分得
\[
V=-\frac{\nu a}{T^n}+\frac{\nu RT}{p}+f(p),
\]
对 $p$ 积分得
\[
V=\frac{\nu RT}{p}+g(T).
\]
合并得
\[
V=C-\frac{\nu a}{T^n}+\frac{\nu RT}{p},
\]
即
\[
p\left(V-C+\frac{\nu a}{T^n}\right)=\nu RT.
\]
\end{solution}
\end{question}

\begin{question}
假定组成一质量为 $M$ 的星体的物质可以近似为单原子分子理想气体，试确定其在过程方程
\[
V=\frac{a_0}{p^{5/6}}
\]
中的比热容，并说明它与宇宙不会陷入热寂的联系。
\begin{solution}
设物质的量为 $\nu=M/\mu$。由 $pV=\nu RT$ 与 $V=a_0p^{-5/6}$，知 $p\propto T^6$、$V\propto T^{-5}$，所以 $dV/dT=-5V/T$。单原子理想气体 $U=3\nu RT/2$，故
\[
C=\frac{\delta Q}{dT}=\frac{dU+p\,dV}{dT}
=\frac32\nu R-5\nu R=-\frac72\nu R.
\]
该过程热容为负，体现了自引力系统的典型负热容：放出能量时反而升温，吸收能量时可能降温。因此引力体系不趋向简单的均匀热平衡，这也是“热寂”图像在有引力系统中不能直接套用的原因之一。
该过程热容为负，体现了自引力系统的典型负热容：放出能量时反而升温，吸收能量时可能降温。因此引力体系不趋向简单的均匀热平衡，这也是"热寂"图像在有引力系统中不能直接套用的原因之一。
\end{solution}
\end{question}

\begin{question}
一级相变的主要特征有哪些？二级相变与一级相变的主要差别有哪些？对有序无序相变，序参量大小和对称性高低有什么关系？超导相变和水的沸腾分别如何用对称性语言描述？
\begin{solution}
一级相变有潜热，熵、体积等热力学势的一阶导数发生跳变，并有相共存和滞后。二级相变无潜热，一阶导数连续，但热容、压缩系数、膨胀系数等二阶导数不连续或发散。高对称相通常序参量为零，低对称相序参量非零。正常金属到超导相是超导序参量出现的对称性破缺过程。水的沸腾是液气一级相变，严格说不属于典型的对称性破缺相变；若按有序度类比，从液体到气体可视作更接近对称性恢复。
\end{solution}
\end{question}

\begin{question}
证明液体的表面内能密度 $u_S$ 可由液体的表面张力系数 $\sigma$ 及温度 $T$ 表示为
\[
u_S=\sigma-T\left(\frac{\partial\sigma}{\partial T}\right)_{A_S},
\]
其中 $A_S$ 为表面面积。
\begin{solution}
表面自由能为
\[
F_S=\sigma A_S.
\]
表面熵为
\[
S_S=-\left(\frac{\partial F_S}{\partial T}\right)_{A_S}
=-A_S\left(\frac{\partial\sigma}{\partial T}\right)_{A_S}.
\]
表面内能 $U_S=F_S+TS_S$，因此
\[
U_S=A_S\left[\sigma-T\left(\frac{\partial\sigma}{\partial T}\right)_{A_S}\right].
\]
两边除以表面积即得
\[
u_S=\sigma-T\left(\frac{\partial\sigma}{\partial T}\right)_{A_S}.
\]
\end{solution}
\end{question}

\textbf{二、计算题（扫描件手写提纲可辨认部分）}

\begin{question}
验证 Maxwell 速率分布。设通过检测装置观测到的速率分布会因粒子通量而按速率加权。

(1) 说明检测到的速率分布应为
\[
F_M(v)=\frac{v f(v)}{\bar v}.
\]

(2) 若检测缝宽或几何线度 $d$ 不可忽略，给出应如何修正测量到的分布。
\begin{solution}
Maxwell 速率分布 $f(v)$ 是容器内分子按速率的数目分布。穿过小孔或到达探测器的分子通量与法向速度成正比；对各向同性气体，固定速率 $v$ 的分子到达探测器的权重与 $v$ 成正比。因此通量加权后的速率分布为
\[
F_M(v)=Cvf(v).
\]
归一化条件 $\int_0^\infty F_M(v)\,dv=1$ 给出 $C=1/\bar v$，故
\[
F_M(v)=\frac{v f(v)}{\bar v}.
\]
若几何线度 $d$ 不可忽略，探测器不会选择单一速率，而是对一个速率窗口取平均。若名义速率为 $V$，由装置几何确定的分辨宽度为 $\Delta V(d)$，则应写为卷积或窗口平均
\[
F_M^{(d)}(V)=\int_0^\infty W_d(V,v)\frac{v f(v)}{\bar v}\,dv,
\qquad
\int_0^\infty W_d(V,v)\,dV=1.
\]
窄缝近似下可用
\[
F_M^{(d)}(V)\simeq\frac1{\Delta V}\int_{V-\Delta V/2}^{V+\Delta V/2}\frac{v f(v)}{\bar v}\,dv,
\]
其中 $\Delta V$ 由飞行距离、转速或缝宽等具体参数决定；当 $d\to0$ 时回到 $F_M(V)=Vf(V)/\bar v$。
\ansmath{F_M(v)=\frac{v f(v)}{\bar v},\qquad F_M^{(d)}=W_d*\left(\frac{v f(v)}{\bar v}\right)}
\end{solution}
\end{question}

\begin{question}
肥皂泡问题。已知表面张力系数 $\sigma$ 和外界压强 $p_0$。

(1) 求半径为 $R$ 时泡内气压 $p$；

(2) 使泡表面带电量 $q$ 后，求新半径 $R'$ 满足的方程；

(3) 说明为什么 $R'>R$。
\begin{solution}
肥皂泡有内、外两个表面，Laplace 压强差为
\[
p-p_0=\frac{4\sigma}{R},
\]
所以
\[
p=p_0+\frac{4\sigma}{R}.
\]
若带电后电荷均匀分布在半径 $R'$ 的球面上，电场产生的向外电压力为
\[
p_e=\frac{\sigma_e^2}{2\varepsilon_0}
=\frac{q^2}{32\pi^2\varepsilon_0 R'^4}.
\]
若泡内气体物质的量不变且过程等温，则
\[
p_{\rm in}'R'^3=pR^3.
\]
力学平衡为“气体压强差 $+$ 电压力 $=$ 表面张力压强”，即
\[
p_{\rm in}'-p_0+p_e=\frac{4\sigma}{R'}.
\]
代入 $p=p_0+4\sigma/R$ 得 $R'$ 的方程
\[
\left(p_0+\frac{4\sigma}{R}\right)\frac{R^3}{R'^3}-p_0
+\frac{q^2}{32\pi^2\varepsilon_0R'^4}
=\frac{4\sigma}{R'}.
\]
带电产生额外向外压力。若仍取 $R'=R$，原来的气体压强差已与表面张力平衡，再加上电压力后向外力过大，故泡会膨胀，平衡半径满足 $R'>R$。
\end{solution}
\end{question}

\begin{question}
地球上有很多雪山，常年积雪不化的最低海拔高度称为雪线高度。
假设常温下水蒸气的凝结热近似为常量
$\lambda_{\mathrm{vp}}=2.4\times10^6\ \mathrm{J/kg}$，
大气中水蒸气的摩尔浓度随温度的变化率近似为常量
$\dfrac{\mathrm{d}c_{\mathrm{vp}}}{\mathrm{d}T}
=3.1\times10^{-4}\ \mathrm{K^{-1}}$，
海平面上大气的温度为 $27^\circ\mathrm{C}$，试确定雪线高度。
\begin{solution}
设近地面大气可看作理想气体。考虑一湿空气微团绝热上升，
微团与外界无热交换，但内部水蒸气可凝结释放潜热。
由热力学第一定律，
\[
\delta Q = C_{p,m}\,\mathrm{d}T - V\,\mathrm{d}p
+ \lambda_{\mathrm{vp}}\,\mathrm{d}m_{\mathrm{vp}} = 0,
\]
其中 $C_{p,m}$ 为空气的定压摩尔热容，$\mathrm{d}m_{\mathrm{vp}}$ 为温度降低
$\mathrm{d}T$ 时微团内凝结的水蒸气质量。

水蒸气的质量可写为 $m_{\mathrm{vp}} = \mu_{\mathrm{vp}} N c_{\mathrm{vp}}$，
其中 $\mu_{\mathrm{vp}}$ 是水蒸气摩尔质量，$N$ 是微团的总物质的量（不变），
$c_{\mathrm{vp}}$ 是水蒸气摩尔浓度。因此
\[
\mathrm{d}m_{\mathrm{vp}} = \mu_{\mathrm{vp}} N\,\mathrm{d}c_{\mathrm{vp}}.
\]

代入第一定律并利用理想气体状态方程 $V = NRT/p$，得
\[
C_{p,m}\,\mathrm{d}T - \frac{NRT}{p}\,\mathrm{d}p
+ \lambda_{\mathrm{vp}}\mu_{\mathrm{vp}} N\,\mathrm{d}c_{\mathrm{vp}} = 0.
\]

除以 $N$（每摩尔形式）：
\[
C_{p,m}\,\mathrm{d}T - \frac{RT}{p}\,\mathrm{d}p
+ \lambda_{\mathrm{vp}}\mu_{\mathrm{vp}}\,\mathrm{d}c_{\mathrm{vp}} = 0.
\]

由静力平衡条件 $\mathrm{d}p = -\rho g\,\mathrm{d}z$ 及理想气体
$\rho = p\bar\mu/(RT)$（$\bar\mu$ 为空气平均摩尔质量），得
\[
\frac{\mathrm{d}p}{\mathrm{d}z} = -\frac{p\bar\mu g}{RT}.
\]

将 $\mathrm{d}p$ 代入，并利用题给 $\mathrm{d}c_{\mathrm{vp}}/\mathrm{d}T$ 为常量，
将 $\mathrm{d}c_{\mathrm{vp}} = \dfrac{\mathrm{d}c_{\mathrm{vp}}}{\mathrm{d}T}\,\mathrm{d}T$ 代入：
\[
C_{p,m}\,\mathrm{d}T
- \frac{RT}{p}\!\left(-\frac{p\bar\mu g}{RT}\,\mathrm{d}z\right)
+ \lambda_{\mathrm{vp}}\mu_{\mathrm{vp}}\frac{\mathrm{d}c_{\mathrm{vp}}}{\mathrm{d}T}\,\mathrm{d}T = 0.
\]

整理得
\[
\left(C_{p,m} + \lambda_{\mathrm{vp}}\mu_{\mathrm{vp}}
\frac{\mathrm{d}c_{\mathrm{vp}}}{\mathrm{d}T}\right)
\frac{\mathrm{d}T}{\mathrm{d}z} = -\bar\mu g,
\]

于是湿绝热直减率为
\[
\frac{\mathrm{d}T}{\mathrm{d}z}
= -\frac{\bar\mu g}
{C_{p,m} + \lambda_{\mathrm{vp}}\mu_{\mathrm{vp}}
\dfrac{\mathrm{d}c_{\mathrm{vp}}}{\mathrm{d}T}}.
\]

从海平面 $z=0$（温度 $T_0$）到雪线高度 $z_{\text{雪线}}$（温度 $T_{\text{雪线}}$），
上式右端各项均为常量，可直接积分：
\[
\int_{T_0}^{T_{\text{雪线}}} \mathrm{d}T
= -\frac{\bar\mu g}
{C_{p,m} + \lambda_{\mathrm{vp}}\mu_{\mathrm{vp}}
\dfrac{\mathrm{d}c_{\mathrm{vp}}}{\mathrm{d}T}}
\int_0^{z_{\text{雪线}}} \mathrm{d}z,
\]

即
\[
T_{\text{雪线}} - T_0
= -\frac{\bar\mu g}
{C_{p,m} + \lambda_{\mathrm{vp}}\mu_{\mathrm{vp}}
\dfrac{\mathrm{d}c_{\mathrm{vp}}}{\mathrm{d}T}}
\,z_{\text{雪线}}.
\]

因此
\[
z_{\text{雪线}}
= \frac{(T_0 - T_{\text{雪线}})
\bigl(
C_{p,m} + \lambda_{\mathrm{vp}}\mu_{\mathrm{vp}}
\dfrac{\mathrm{d}c_{\mathrm{vp}}}{\mathrm{d}T}
\bigr)}
{\bar\mu g}.
\]

海平面温度 $T_0 = 27^\circ\mathrm{C} = 300.15\ \mathrm{K}$，
雪线温度即冰点 $T_{\text{雪线}} = 0^\circ\mathrm{C} = 273.15\ \mathrm{K}$。
代入数据
\[
\begin{aligned}
C_{p,m} &= \frac{7}{2}R = 29.085\ \mathrm{J/(mol\cdot K)},\\[2pt]
\bar\mu &= 29\times10^{-3}\ \mathrm{kg/mol},\\[2pt]
\mu_{\mathrm{vp}} &= 18\times10^{-3}\ \mathrm{kg/mol},\\[2pt]
\frac{\mathrm{d}c_{\mathrm{vp}}}{\mathrm{d}T}
        &= 3.1\times10^{-4}\ \mathrm{K^{-1}},\\[2pt]
g &= 9.8\ \mathrm{m/s^2},
\end{aligned}
\]

得
\[
z_{\text{雪线}} \approx 4.0\times10^3\ \mathrm{m}.
\]
\end{solution}
\end{question}


\begin{question}
二级相变中，证明 Ehrenfest 关系
\[
\left(\frac{dp}{dT}\right)_{\rm tr}=\frac{\Delta\alpha}{\Delta\kappa_T}.
\]
\begin{solution}
二级相变线上两相的体积相等且连续：$V_1=V_2$。沿相变线取微分，有
\[
dV_1=dV_2.
\]
对任一相
\[
dV=V\alpha\,dT-V\kappa_T\,dp.
\]
二级相变处 $V_1=V_2=V$，代入得
\[
V\alpha_1dT-V\kappa_{T1}dp
=V\alpha_2dT-V\kappa_{T2}dp.
\]
整理为
\[
(\alpha_2-\alpha_1)dT=(\kappa_{T2}-\kappa_{T1})dp.
\]
因此
\[
\left(\frac{dp}{dT}\right)_{\rm tr}
=\frac{\alpha_2-\alpha_1}{\kappa_{T2}-\kappa_{T1}}
=\frac{\Delta\alpha}{\Delta\kappa_T}.
\]
\end{solution}
\end{question}

